\begin{exo}
	Dans le tableau ci-dessous, on note :
	\begin{itemize}
		\item $F$ : \og{} La personne est une fille \fg{}
		\item $V$ : \og{} La personne a plus de 20 ans \fg{}
	\end{itemize}

	On choisit une personne au hasard dans ce groupe.

	\begin{center}
		\renewcommand{\arraystretch}{1.2}
		\begin{tabular}{|c|C{2cm}|C{2cm}|c|}
			\hline
			               & $F$ & $\overline{F}$ & \textbf{Total} \\ \hline
			$V$            & 6   & 44             & 50             \\ \hline
			$\overline{V}$ & 56  & 59             & 115            \\ \hline
			\textbf{Total} & 62  & 103            & 165            \\ \hline
		\end{tabular}
	\end{center}

	À l'aide du tableau, déterminer les probabilités suivantes (on donnera les résultats sous forme de fraction) :

	\begin{enumerate}
		\item \textbf{Proba. marginales :} $P\pa{F}$ et $P\pa{V}$
		\item \textbf{Proba. d'intersection :} $P\pa{F \cap V}$ et $P\pa{\overline{F} \cap \overline{V}}$
		\item \textbf{Proba. conditionnelles :} $P_{V}\pa{\overline{F}}$ et $P_{F}\pa{V}$
	\end{enumerate}
\end{exo}
